\chapter{Conclusions}

\section{Summary}
This thesis studied how to detect very small human targets in aerial disaster
imagery with a model small enough to deploy, and it did so in two phases. In the
first, an additive ablation on a YOLO11m baseline measured the effect of CBAM
attention, a high-resolution P2 detection scale, and a state-space neck on the
C2A dataset under one training and evaluation protocol. The P2 scale was the
component that mattered: it raised very-tiny recall and $\mathrm{AP}_{50}$ at a
small parameter and latency cost. CBAM was near-free and slightly improved
operational scores while reducing model size. The state-space neck added
parameters and roughly tripled latency without improving accuracy, and was
excluded. The recommended CBAM~+~P2 model is competitive with much larger
published detectors on the same data at about a third of their parameter count.
The second phase hardened and extended that result. A split audit showed the
official benchmark leaks scenes between training and test sets at a measured
cost of $1.6$ points of $\mathrm{AP}_{50}$, and all further work moved to a
scene-disjoint protocol. Under it, a blend of a normalised Wasserstein
similarity into the task-aligned assigner raised sub-8-pixel recall by $2.1$
points with a bootstrap interval that excludes zero, at a measured cost on
larger targets; a frequency-gated architecture branch failed the same paired
test and joined the state-space neck as a documented negative. Real footage
collected and annotated for this thesis then showed that the synthetic-trained
model transfers poorly zero-shot ($\mathrm{AP}_{50}$ $0.335$ on drone imagery,
$0.133$ on real disaster frames), that a labelled budget of a few hundred frames
closes most of the drone gap ($0.795$) and triples the within-video disaster
score ($0.411$) at almost no cost to the source benchmark, and that this
adaptation does not carry to a disaster event held out from training, where the
gain over the un-adapted model is not statistically significant. Two findings
organise the whole: on this task, training supervision beat added architecture
everywhere it was tried, and adaptation gains are specific to the footage they
come from.

\section{Limitations}
The study has clear limits. The primary benchmark, C2A, is semi-synthetic, and
its paste-based annotations put a ceiling on localisation accuracy that all
compared models share. Every headline training run uses a single seed; the
bootstrap intervals of Chapter~IV treat test-set variation only, and the
multi-seed replication of the assignment pair was still in progress at the time
of writing. The real evaluation sets are small (23 to 60 frames each), so their
scores carry wide intervals and should be read as evidence of direction rather
than precise levels. The within-video disaster result shares source videos
between the training and evaluation frames, deliberately and with de-duplication,
and is reported as same-event adaptation, not generalisation; the honest
cross-event number is the one from the held-out event. Recall of sub-8-pixel
people in real disaster footage is zero for every model tested here, so the
regime the thesis targets on synthetic data remains open on real data. The
adaptation fine-tune also cost 13 points of precision on the held-out event, so
recall-side gains were partly paid for in false alarms. On-device measurement on
drone hardware is still not covered.

\section{Recommendations and Future Works}
The limits above set the agenda. The first need is scale on the real-disaster
side: the cross-event result shows that 60 frames from two events adapt the
model to those events only, so a multi-event disaster training set, on the order
of many events rather than many frames per event, is the most direct route to
generalisation, and the annotation pipeline built here (frame extraction,
perceptual-hash de-duplication, assisted labelling, automated audits) makes each
additional event cheap. The multi-seed replication of the assignment pair should
be completed and folded into the reported table. The sub-8-pixel failure on real
footage suggests combining the assignment blend with higher-resolution real
training data, or with the slicing of Chapter~IV at inference time, and
measuring whether the two compose. Cross-dataset validation on real corpora such
as SARD~\cite{sambolek2021sard}, HERIDAL~\cite{bozic2019heridal}, and
VisDrone~\cite{zhu2021visdrone} remains worthwhile. On the deployment side, the
lightweight variant should be exported and measured on embedded hardware. The
architecture-side negatives do not close their directions: a state-space or
frequency-domain design built specifically around the geometry of sub-8-pixel
targets, rather than adapted from generic modules, remains untested, and the
newest YOLO generation should be re-run through the same protocol as a stronger
starting point.
